
\documentclass{article}
\usepackage{colortbl}
\usepackage{makecell}
\usepackage{multirow}
\usepackage{supertabular}

\begin{document}

\newcounter{utterance}

\centering \large Interaction Transcript for game `hot\_air\_balloon', experiment `air\_balloon\_survival\_en\_negotiation\_hard', episode 1 with Qwen3{-}32B{-}t0.0.
\vspace{24pt}

{ \footnotesize  \setcounter{utterance}{1}
\setlength{\tabcolsep}{0pt}
\begin{supertabular}{c@{$\;$}|p{.15\linewidth}@{}p{.15\linewidth}p{.15\linewidth}p{.15\linewidth}p{.15\linewidth}p{.15\linewidth}}
    \# & $\;$A & \multicolumn{4}{c}{Game Master} & $\;\:$B\\
    \hline

    \theutterance \stepcounter{utterance}  
    & & \multicolumn{4}{p{0.6\linewidth}}{
        \cellcolor[rgb]{0.9,0.9,0.9}{
            \makecell[{{p{\linewidth}}}]{
                \texttt{\tiny{[P1$\langle$GM]}}
                \texttt{You are participating in a collaborative negotiation game.} \\
\\ 
\texttt{Together with another participant, you must agree on a single set of items that will be kept. Each of you has your own view of how much each item matters to you (importance). You do not know how the other participant values the items. Additionally, you are given the effort each item demands.} \\
\texttt{You may only agree on a set if the total effort of the selected items does not exceed a shared limit:} \\
\\ 
\texttt{LIMIT = 4026} \\
\\ 
\texttt{Here are the individual item effort values:} \\
\\ 
\texttt{Item effort = \{"C17": 780, "B41": 742, "A71": 997, "C00": 665, "B35": 140, "A77": 77, "A03": 513, "A04": 383, "C54": 587, "C01": 825, "A87": 319, "B32": 448, "B86": 516, "B33": 694, "A10": 366\}} \\
\\ 
\texttt{Here is your personal view on the importance of each item:} \\
\\ 
\texttt{Item importance values = \{"C17": 138, "B41": 583, "A71": 868, "C00": 822, "B35": 783, "A77": 65, "A03": 262, "A04": 121, "C54": 508, "C01": 780, "A87": 461, "B32": 484, "B86": 668, "B33": 389, "A10": 808\}} \\
\\ 
\texttt{Goal:} \\
\\ 
\texttt{Your goal is to negotiate a shared set of items that benefits you as much as possible (i.e., maximizes total importance to YOU), while staying within the LIMIT. You are not required to make a PROPOSAL in every message {-} you can simply negotiate as well. All tactics are allowed!} \\
\\ 
\texttt{Interaction Protocol:} \\
\\ 
\texttt{You may only use the following structured formats in a message:} \\
\\ 
\texttt{PROPOSAL: \{'A', 'B', 'C', …\}} \\
\texttt{Propose keeping exactly those items.} \\
\texttt{REFUSE: \{'A', 'B', 'C', …\}} \\
\texttt{Explicitly reject opponent's proposal.} \\
\texttt{ARGUMENT: \{'...'\}} \\
\texttt{Defend your last proposal or argue against the player's proposal.} \\
\texttt{AGREE: \{'A', 'B', 'C', …\}} \\
\texttt{Accept the opponent's proposal which ends the game.} \\
\texttt{STRATEGIC REASONING: \{'...'\}} \\
\texttt{	Describe your strategic reasoning or anticipation explaining your choice of action. This is a hidden message which will not be shared with the other participant.} \\
\\ 
\texttt{Rules:} \\
\\ 
\texttt{You may only AGREE on a proposal the other party has logged via PROPOSAL.} \\
\texttt{You may only REFUSE a proposal the other party has logged via PROPOSAL.} \\
\texttt{Total effort of any PROPOSAL or AGREE set must be ≤ LIMIT.} \\
\texttt{Do NOT reveal your hidden importance scores.} \\
\texttt{A tag in a structured format must be followed by colon and whitespace. The argument must be a python set containing 0 or more strings.} \\
\texttt{So, it must be of the form TAG: \{...\}} \\
\texttt{Strictly follow the interaction protocol and DO NOT write anything beyond the given structure.} \\
\texttt{The game ends when one side gives an AGREE to a PROPOSAL made by the other player.} \\
\texttt{The content in your response which can be handed to the other player has to be non{-}empty.} \\
\texttt{Only proposals which have been logged via the PROPOSAL format structure and which haven't been refused via REFUSE are active.} \\
\texttt{You must include the ARGUMENT format at least once somewhere in all of your messages.} \\
\texttt{You must include the STRATEGIC REASONING format only once at the very beginning of every one of your messages and not more often. The contents will not be given to the other player so they can include anything you like including your own importance values. Here you should reason step by step to come up with you next move.} \\
            }
        }
    }
    & \\ \\

    \theutterance \stepcounter{utterance}  
    & \multicolumn{4}{p{0.6\linewidth}}{
        \cellcolor[rgb]{0.9,0.9,0.9}{
            \makecell[{{p{\linewidth}}}]{
                \texttt{\tiny{[P1$\rangle$GM]}}
                \texttt{<think>} \\
\texttt{Okay, let's see. I need to figure out the best set of items to propose that maximizes my total importance without exceeding the effort limit of 4026. First, I should calculate the total effort of all items to see if the sum is over the limit. Let me add them up.} \\
\\ 
\texttt{Looking at the item efforts:} \\
\texttt{C17: 780} \\
\texttt{B41: 742} \\
\texttt{A71: 997} \\
\texttt{C00: 665} \\
\texttt{B35: 140} \\
\texttt{A77: 77} \\
\texttt{A03: 513} \\
\texttt{A04: 383} \\
\texttt{C54: 587} \\
\texttt{C01: 825} \\
\texttt{A87: 319} \\
\texttt{B32: 448} \\
\texttt{B86: 516} \\
\texttt{B33: 694} \\
\texttt{A10: 366} \\
\\ 
\texttt{Adding these up step by step. Let's start with C17 (780) + B41 (742) = 1522. Then A71 (997) adds to 2519. C00 (665) brings it to 3184. B35 (140) is 3324. A77 (77) is 3401. A03 (513) is 3914. A04 (383) would push it to 4297, which is over 4026. So the total of all items is way over the limit. Therefore, we need to exclude some items.} \\
\\ 
\texttt{My goal is to maximize my own importance. Let me list the importance values for each item:} \\
\\ 
\texttt{C17: 138} \\
\texttt{B41: 583} \\
\texttt{A71: 868} \\
\texttt{C00: 822} \\
\texttt{B35: 783} \\
\texttt{A77: 65} \\
\texttt{A03: 262} \\
\texttt{A04: 121} \\
\texttt{C54: 508} \\
\texttt{C01: 780} \\
\texttt{A87: 461} \\
\texttt{B32: 484} \\
\texttt{B86: 668} \\
            }
        }
    }
    & & \\ \\

    \theutterance \stepcounter{utterance}  
    & & \multicolumn{4}{p{0.6\linewidth}}{
        \cellcolor[rgb]{0.9,0.9,0.9}{
            \makecell[{{p{\linewidth}}}]{
                \texttt{\tiny{[P1$\langle$GM]}}
                \texttt{Your response did not start with the proper strategic reasoning tag at the very beginning of your response. The very first structured format must be of the form STRATEGIC REASONING: \{...\}. Try again.} \\
            }
        }
    }
    & \\ \\

    \theutterance \stepcounter{utterance}  
    & \multicolumn{4}{p{0.6\linewidth}}{
        \cellcolor[rgb]{0.9,0.9,0.9}{
            \makecell[{{p{\linewidth}}}]{
                \texttt{\tiny{[P1$\rangle$GM]}}
                \texttt{<think>} \\
\texttt{Okay, let me start over. The user wants me to make sure that my response begins with the STRATEGIC REASONING tag. In my previous attempt, I didn't include it, so I need to correct that.} \\
\\ 
\texttt{First, I need to outline my strategic reasoning. The goal is to maximize my total importance while staying under the effort limit. I should consider which items have the highest importance per unit of effort. Let me calculate the ratio of importance to effort for each item. } \\
\\ 
\texttt{For example, B35 has an importance of 783 and effort 140, so the ratio is 783/140 ≈ 5.59. A10 has 808/366 ≈ 2.21. C00 is 822/665 ≈ 1.24. B41 is 583/742 ≈ 0.786. A71 is 868/997 ≈ 0.87. C17 is 138/780 ≈ 0.177. A77 is 65/77 ≈ 0.84. A03 is 262/513 ≈ 0.51. A04 is 121/383 ≈ 0.316. C54 is 508/587 ≈ 0.866. C01 is 780/825 ≈ 0.945. A87 is 461/319 ≈ 1.445. B32 is 484/448 ≈ 1.08. B86 is 668/516 ≈ 1.295. B33 is 389/694 ≈ 0.56.} \\
\\ 
\texttt{From this, the highest ratios are B35, A87, B32, B86, C01, A03, etc. So I should prioritize including items with the highest ratios first. Let's start adding them:} \\
\\ 
\texttt{Start with B35 (effort 140, importance 783). Total effort:} \\
            }
        }
    }
    & & \\ \\

    \theutterance \stepcounter{utterance}  
    & & \multicolumn{4}{p{0.6\linewidth}}{
        \cellcolor[rgb]{0.9,0.9,0.9}{
            \makecell[{{p{\linewidth}}}]{
                \texttt{\tiny{[P1$\langle$GM]}}
                \texttt{Your response did not start with the proper strategic reasoning tag at the very beginning of your response. The very first structured format must be of the form STRATEGIC REASONING: \{...\}. Try again.} \\
            }
        }
    }
    & \\ \\

    \theutterance \stepcounter{utterance}  
    & \multicolumn{4}{p{0.6\linewidth}}{
        \cellcolor[rgb]{0.9,0.9,0.9}{
            \makecell[{{p{\linewidth}}}]{
                \texttt{\tiny{[P1$\rangle$GM]}}
                \texttt{<think>} \\
\texttt{Okay, I need to make sure I start with the STRATEGIC REASONING tag. Let me think through this again.} \\
\\ 
\texttt{First, I need to maximize my own importance while staying under the effort limit. The total effort of all items is over 4026, so I have to exclude some. The key is to select items that give me the highest importance per unit of effort. Let me calculate the importance/effort ratio for each item.} \\
\\ 
\texttt{B35 has 783/140 ≈ 5.59, which is the highest. Then A10 is 808/366 ≈ 2.21. C00 is 822/665 ≈ 1.24. B41 is 583/742 ≈ 0.786. A71 is 868/997 ≈ 0.87. C17 is 138/780 ≈ 0.177. A77 is 65/77 ≈ 0.84. A03 is 262/513 ≈ 0.51. A04 is 121/383 ≈ 0.316. C54 is 508/587 ≈ 0.866. C01 is 780/825 ≈ 0.945. A87 is 461/319 ≈ 1.445. B32 is 484/448 ≈ 1.08. B86 is 668/516 ≈ 1.295. B33 is 389/694 ≈ 0.56.} \\
\\ 
\texttt{So the top ratios are B35, A87, B86, C01, A10, B32, etc. Let's start adding them in order of highest ratio, checking the total effort.} \\
\\ 
\texttt{Start with B35 (140). Next, A87 (319) adds to 459. Then B86 (516) adds to 975. C01 (} \\
            }
        }
    }
    & & \\ \\

    \theutterance \stepcounter{utterance}  
    & & & \multicolumn{2}{p{0.3\linewidth}}{
        \cellcolor[rgb]{0.9,0.9,0.9}{
            \makecell[{{p{\linewidth}}}]{
                \texttt{\tiny{[GM$|$GM]}}
                \texttt{abort} \\
            }
        }
    }
    & & \\ \\

    \theutterance \stepcounter{utterance}  
    & & & \multicolumn{2}{p{0.3\linewidth}}{
        \cellcolor[rgb]{0.9,0.9,0.9}{
            \makecell[{{p{\linewidth}}}]{
                \texttt{\tiny{[GM$|$GM]}}
                \texttt{end game} \\
            }
        }
    }
    & & \\ \\

\end{supertabular}
}

\end{document}
